\section{Experiments}
\subsection{Setup}
\paragraph{Implementation Details}
The training is conducted with an effective batch size of $128$ over $50k$ steps on a $7,600$-hour audio-visual datasets built by our data curation pipeline,
using the AdamW optimizer with a learning rate of $5e-6$. We employ Fully Sharded Data Parallel (FSDP) for distributed training across multiple nodes, with a gradient accumulation step of $4$.
\paragraph{Compared Methods} We conduct a comprehensive evaluation of our model by benchmarking it against a suite of state-of-the-art baselines across several distinct generation tasks:
\begin{itemize}[leftmargin=*]
    \item Video Generation: We compare our model against leading text-to-video systems, including Wan2.2 (14B)~\cite{wan2025wan}, HunyuanVideo~\cite{kong2024hunyuanvideo}, CogVideoX-1.5, (5B)~\cite{yang2024cogvideox}, Kling 2.1~\cite{kuaishou2024kling}, and SeeDance 1.0~\cite{gao2025seedance}.
    \item Audio Generation: For the audio modality, we benchmark against established text-to-audio models such as Stable Audio Open~\cite{evans2025stable} and AudioLDM2~\cite{audioldm2-2024taslp}.
    \item Text-to-Speech (TTS): As a supplementary evaluation of vocal synthesis, we compare our model's performance against specialized TTS models, including CosyVoice~\cite{du2024cosyvoice}, CosyVoice2~\cite{du2024cosyvoice2}, and VibeVoice~\cite{peng2025vibevoice}.
    \item Audio-to-Video methods: We also compare out model with talking-based audio to video methods such as FantastyTalking~\cite{wang2025fantasytalking} and Wan-S2V~\cite{gao2025wan}.
    \item Video-to-Audio methods: As a closely related task to joint audio-visual generation, we also benchmarked our model on the video-to-audio (V2A) task, drawing comparisons with state-of-the-art methods such as HunyuanVideo-Foley~\cite{shan2025hunyuanvideo}.
    \item Joint Audio-Video Generation: For our core task of synchronous audio-visual synthesis, we compare our method against existing prompt-based joint generation models: SVG~\cite{ishii2024simple}. Since methods such as MM-Diffusion~\cite{ruan2023mm} and R-FLAV~\cite{ergasti2025r} are class-conditional generative models, a direct comparison with our approach is not applicable, we only compare with SVG in our report.
\end{itemize}
It is important to note that our model is constructed by stitching the pre-trained WAN2.1 (1.3B) and Step-Ace (3.5B) models. Consequently, the comparisons against the aforementioned state-of-the-art baselines are intended to provide a qualitative reference and situate our work, rather than to make a direct claim of superior performance.

\paragraph{Benchmark} To construct our evaluation set, we curated $600$ image-text prompt pairs from a multitude of sources. These sources encompass frames extracted from YouTube videos, BiliBili videos, TikTok clips, movies, and anime; images generated by AI models~\cite{jimeng,wu2025qwen}; and a collection of images from public websites. Our dataset comprises three subsets.
\begin{itemize}[leftmargin=*]
    \item Set1-I contains image-text pairs (including AI-generated, web-crawled, and media screenshots), for which video/audio captions and speech content were produced using LLMs~\cite{xu2025qwen2} and manual annotation, comprising a total of $205$ samples. Statistical results in~\ref{fig:unibench_statistics_set1}.
    \item Set2-V consists of video clips from YouTube and Bilibili, which were annotated with LLM-generated~\cite{xu2025qwen2} captions and Whisper-based ASR~\cite{radford2023robust} transcripts, followed by human verification, comprising a total of $295$ samples. Statistical results in~\ref{fig:unibench_statistics_set2}.
    \item Set3-Ted (the Verse-Ted subset) includes TED Talks from September 2025, processed with the same annotation pipeline as Set2, comprising a total of $100$ samples. 
\end{itemize}
Our collected test data is highly diverse, encompassing a wide spectrum of audio categories: human speech; animal vocalizations (e.g., bird chirping, cat meowing); instrumental music (e.g., piano, guitar); natural sounds (e.g., thunder, rain); human-object interactions (e.g., keyboard typing, cooking, chopping vegetables); object-object interactions (e.g., glass shattering, marbles dropping); mechanical sounds (e.g., trains, airplanes), and so on. The complete statistical results of set1 and set2 are shown in~\ref{fig:unibench_statistics_all}.

\begin{figure*}
    \centering
    \begin{subfigure}[h]{0.3\textwidth}
    \centering
    \includegraphics[width=1.0\linewidth]{figures/audio_statistics_all.pdf}
    % \vspace{-3mm}
    \caption{Statistical results of set1/2.}
    \label{fig:unibench_statistics_all}
    % \vspace{-5mm}
    \end{subfigure}
    \hspace{0.01\textwidth} % 可选的右侧间距
    \vrule
    \hspace{0.01\textwidth} % 可选的左侧间距
    \begin{subfigure}[h]{0.3\textwidth}
    \centering
    \includegraphics[width=1.0\linewidth]{figures/audio_statistics_set1.pdf}
    % \vspace{-3mm}
    \caption{Statistical results of set1.}
    \label{fig:unibench_statistics_set1}
    % \vspace{-5mm}
    \end{subfigure}
    \hspace{0.01\textwidth} % 可选的右侧间距
    \vrule
    \hspace{0.01\textwidth} % 可选的左侧间距
    \begin{subfigure}[h]{0.3\textwidth}
    \centering
    \includegraphics[width=1.0\textwidth]{figures/audio_statistics_set2.pdf}
    % \vspace{-3mm}
    \caption{Statistical results of set2.}
    \label{fig:unibench_statistics_set2}
    \end{subfigure}
    
    \caption{Statistical results of \textbf{Verse-Bench}. Best viewed with \textbf{zoom-in}. A larger, high-resolution version is available in Appendix~\ref{lab:verse-bench-fig}}
    % \vspace{-7mm}
\end{figure*}

\paragraph{Evaluation Protocol}
We quantitatively evaluate our method against baselines on the Verse-Bench benchmark. The evaluation is structured across $6$ distinct generation tasks, each with a tailored set of metrics.

\begin{itemize}[leftmargin=*]
    \item Video Generation: Performance is assessed on three criteria:
    \begin{itemize}
        \item Motion Score (MS): This metric quantifies motion dynamics, calculated from the normalized optical flow magnitude detected by the RAFT model~\cite{teed2020raft}.
        \item Aesthetic Score (AS): This is a composite score averaging three components: fidelity, measured by MANIQA~\cite{yang2022maniqa} to penalize blur and artifacts, and aesthetic quality, evaluated by both aesthetic-predictor-v2-5~\cite{aesthetic-predictor-v2-5} and Musiq~\cite{ke2021musiq}.
        \item ID Consistency (ID): To measure identity preservation, we compute the mean DINOV3~\cite{simeoni2025dinov3} feature similarity between the reference image and each generated frame.
    \end{itemize} 
        
    \item Audio Generation: We evaluate audio quality from three perspectives:
    \begin{itemize}
        \item Distributional Similarity: We measure the Fréchet Distance (FD) and Kullback-Leibler (KL) divergence between the generated and real data distributions, using features extracted from PANNs~\cite{kong2020panns} and PaSST~\cite{koutini2021efficient}.
        \item Semantic Consistency: The alignment between the audio and the input text is measured by the LAION-CLAP~\cite{wu2023large} score.
        \item Quality and Diversity: We report the Inception Score (IS) calculated with a PANNs classifier. Additionally, we use AudioBox-Aesthetics~\cite{tjandra2025meta} to assess Production Quality (PQ), Production Complexity (PC), Content Enjoyment (CE), and Content Usefulness (CU).
    \end{itemize}
    \item Text-to-Speech (TTS): We evaluate synthesis accuracy using the Word Error Rate (WER), which is derived by transcribing the generated audio with the Whisper-large-v3 model~\cite{radford2023robust}.
    \item Audio-to-Video: We evaluate this task using the same criteria as the video generation task, additionally providing a SyncNet~\cite{chung2016out} confidence score to assess lip-sync accuracy.
    \item Video-to-Audio: This task use all metrics from audio generation tasks. Furthermore, we introduce the Audio-Video Alignment (AV-A) metric to specifically quantify the temporal synchronization between the generated audio and video streams, which is computed via Synchformer~\cite{iashin2024synchformer}.
    \item Joint Audio-Video Generation: For this task, we use all relevant metrics from the individual tasks above. 
\end{itemize}

The evaluation of our models and their components is conducted across the three test sets as follows:
\begin{itemize}[leftmargin=*]
    \item The video generation models and SVG are evaluated on Set 1 and Set 2.
    \item For the audio generation model is evaluated on Set1 and Set2.
    \item The Text-to-Speech (TTS) model is primarily evaluated on Set 3.
    \item The Audio-to-Video (A2V) model is evaluated exclusively on Set 3, while the Video-to-Audio (V2A) model is evaluated exclusively on Set 2.
    \item Finally, our complete Universe-1 model is benchmarked against all three test sets.
\end{itemize}
For audio generation, we evaluate the metrics CE, CU, PC, and PQ on Set 1. On Set 2, the evaluation is expanded to include FD, KL, and CS in addition to the aforementioned metrics. Furthermore, LSE-C is evaluated exclusively on Set 3, while the AV-A metric is also applied to Set 1 when evaluating UniVerse-1 and SVG.

\subsection{Quantitative Evaluation}
% We compared Universe-1 with various state-of-the-art (SOTA) models as shown in Tab.~\ref{tab:uni_bench}. In video generation, our model surpasses all competing methods in identity preservation (ID: 0.89) and achieves a competitive score in action similarity (AS: 0.47). For audio generation, while not outperforming specialized audio models on all metrics, it shows strong performance in pitch correlation (PC: 2.49), surpassing most baselines. The most significant advantage of our model is demonstrated in audio-video coherence tasks, where its lip-sync score (LSE-C: 7.10) is substantially higher than other methods, strongly validating the effectiveness of our joint generation framework. Finally, in the TTS task, the Word Error Rate (WER) of 0.18 indicates performance comparable to that of dedicated TTS models.our methods. 

We compare Universe-1 against a range of state-of-the-art (SOTA) models specialized for specific tasks as shown in Tab.~\ref{tab:uni_bench}. As a unified joint generation model, a direct comparison of single-modality metrics against these "expert" models presents inherent complexities. Nevertheless, our model demonstrates robust capabilities across multiple dimensions.
In terms of video quality, Universe-1 achieves the highest score in identity preservation (ID: 0.89), showcasing its superior ability to maintain subject consistency throughout the generation process. For audio quality, while there is a gap compared to leading audio-only generation models, our model obtains a highly competitive score in pitch correlation (PC: 2.49).
The core strength of our model lies in synchronous audio-video generation. It is crucial to interpret the metrics for such joint generation tasks with caution. For instance, in the video-to-audio (V2A) setting, the AV-A metric for a model like Hunyuanvideo-Foley is calculated with a ground-truth video, whereas our model generates both modalities simultaneously. Therefore, AV-A must be considered in conjunction with the audio-text CLAP score (CS) for a holistic assessment. From this integrated perspective, our model (AV-A: 0.23, CS: 0.16) demonstrates a better overall audio-visual content consistency than SVG (AV-A: 0.09, CS: 0.08).
Similarly, for the lip-sync (LSE-C) metric, our model's score (1.34) is evaluated on fully generated audio and video, making it susceptible to the quality of both generated modalities. In contrast, audio-to-video (A2V) methods like Wan-S2V are evaluated using ground-truth audio, which naturally yields a higher score (6.49). Despite this evaluation disparity, Universe-1 achieves promising results as the first open-source joint generation framework of its kind, establishing a solid foundation for future research.

\begin{table}
\vspace{-1em}
\small
\centering
\resizebox{\linewidth}{!}{
    \renewcommand{\arraystretch}{1.8}
    \begin{tabular}{l|c|c|c|c|c|c|c|c|c|c|c|c|c|c}
    \toprule
    % --- 表头部分 ---
    % 第一行主分类表头，使用 \multicolumn 合并单元格
    \multicolumn{2}{c|}{} & \multicolumn{3}{c|}{video} & \multicolumn{7}{c|}{audio} & \multicolumn{1}{c|}{tts} & \multicolumn{2}{c}{Audio-Video} \\
    \midrule
    % 第二行子分类表头
     & method & MS$~\uparrow$ & AS$~\uparrow$ & ID$~\uparrow$ & FD$~\downarrow$ & KL$~\downarrow$ & CS$~\uparrow$ & CE$~\uparrow$ & CU$~\uparrow$ & PC$~\downarrow$ & PQ$~\uparrow$ & WER$~\downarrow$ & LSE-C$~\uparrow$ & AV-A$~\downarrow$ \\
    \midrule
    
    % --- 数据行部分 ---
    % Video 部分，跨越5行
    \multirow{5}{*}{\makecell{video \\ methods}} 
     & CogVideox1.5 5B~\cite{yang2024cogvideox}         & 0.43 & 0.44 & 0.83 & - & - & - & - & - & - & - & - & - & - \\ \cline{2-15} 
     & Wan2.2-14B~\cite{wan2025wan}                     & \textbf{1.04} & \textbf{0.50} & \underline{0.88} & - & - & - & - & - & - & - & - & - & - \\ \cline{2-15} 
     & Kling2.1~\cite{kuaishou2024kling}                & 0.31 & 0.41 & 0.85 & - & - & - & - & - & - & - & - & - & - \\ \cline{2-15}
     & SeeDance1.0~\cite{gao2025seedance}               & \underline{0.50} & \underline{0.47} & 0.86 & - & - & - & - & - & - & - & - & - & - \\
    \midrule
    % Audio 部分，跨越2行
    \multirow{2}{*}{\makecell{audio \\ methods}} 
     & Stable Audio~\cite{evans2025stable}              & - & - & - & \underline{1.13} & \underline{1.38} & \underline{0.28} & 3.88 & \textbf{6.15} & \underline{2.60} & 6.53 & - & - & - \\ \cline{2-15} 
     & AudioLdm2~\cite{audioldm2-2024taslp}             & - & - & - & 1.21 & 2.30 & 0.24 & \underline{3.98} & \underline{5.88} & 3.43 & 6.04 & - & - & - \\
    \midrule
    % TTS 部分，跨越3行
    \multirow{3}{*}{\makecell{tts \\ methods}}   
     & Cosyvoice~\cite{du2024cosyvoice}                 & - & - & - & - & - & - & - & - & - & - & 0.17 & - & - \\ \cline{2-15}  
     & Cosyvoice2~\cite{du2024cosyvoice2}               & - & - & - & - & - & - & - & - & - & - & \underline{0.16} & - & - \\ \cline{2-15}
     & Vibevoice~\cite{peng2025vibevoice}               & - & - & - & - & - & - & - & - & - & - & \textbf{0.15} & - & - \\
    \midrule
    % Audio2Video 部分，跨越2行
    \multirow{2}{*}{\makecell{a2v \\ methods}}   
     & FantastyTalking~\cite{wang2025fantasytalking}    & 0.07 & 0.42 & 0.87 & - & - & - & - & - & - & - & - & 2.68 & - \\ \cline{2-15}  
     & Wan-S2V~\cite{gao2025wan}                        & 0.17 & 0.46 & \textbf{0.89} & - & - & - & - & - & - & - & - & 6.49 & - \\
    \midrule
    % Video2Audio 部分，跨越1行
    \multirow{1}{*}{\makecell{v2a \\ methods}}     
     & Hunyuanvideo-Foley~\cite{shan2025hunyuanvideo}   & - & - & - & \textbf{0.82} & \textbf{1.27} & \textbf{0.40} & \textbf{4.04} & 5.72 & 3.09 & 6.27 & - & - & 0.78 \\
    \midrule
    % AV joint 部分，跨越2行
    \multirow{2}{*}{\makecell{joint \\ methods}}  
     & SVG~\cite{ishii2024simple}                       & 0.40 & 0.41 & 0.25 & 1.55 & 3.62 & 0.08 & 2.93 & 5.50 & \textbf{2.35} & 6.26 & - & - & 0.09 \\ \cline{2-15} 
     & \colorbox{CornflowerBlue}{\textbf{Ours}}         & 0.20 & 0.47 & 0.89 & 1.25 & 2.70 & 0.16 & 3.53 & 4.61 & 2.49 & 5.20 & 0.18 & 1.34 & 0.23 \\ 
    \bottomrule
    \end{tabular}
}
\vspace{1em}
\caption{Quantitative results compared with baselines on Verse-Bench.}
\label{tab:uni_bench}
\vspace{-1em}
\end{table}


% \subsection{Qualitative Evaluation}
% \subsection{User Study}
\subsection{Ablation Study}
We performed an ablation study to investigate the contributions of our Low Quality data Loss Strategy(LQLS) and Independent Noise Sampling Strategy(INSS), as shown in Tab.~\ref{tab:ablation}. The findings indicate that LQLS provides improvement accross video quality and consistency ID, thus confirming its efficacy and positive impact on training. Furthermore, the results for INSS demonstrate a significant enhancement in audio generation quality, validating the effectiveness of this approach.

\begin{table}
% \vspace{-1em}
\small
\centering
\resizebox{\linewidth}{!}{
    \renewcommand{\arraystretch}{1.8}
    \begin{tabular}{l|c|c|c|c|c|c|c|c|c|c|c|c|c}
    \toprule
    % --- 表头部分 ---
    method & MS$~\uparrow$ & AS$~\uparrow$ & ID$~\uparrow$ & FD$~\downarrow$ & KL$~\downarrow$ & CS$~\uparrow$ & CE$~\uparrow$ & CU$~\uparrow$ & PC$~\downarrow$ & PQ$~\uparrow$ & WER$~\downarrow$ & LSE-C$~\uparrow$ & AV-A$~\downarrow$ \\
    \midrule
    % --- 数据行部分 ---
     w/o LQLS                                  & 0.38 & 0.44 & 0.78 & 1.26 & 2.84 & 0.15 & 3.38 & 4.35 & 2.58 & 4.97 & 0.16 & 1.35 & 0.28 \\ \cline{1-14}
     w/o INSS                                  & 1.10 & 0.43 & 0.75 & 1.43 & 3.51 & 0.11 & 2.44 & 3.14 & 2.92 & 3.99 & 0.38 & 0.99 & 0.18 \\ \cline{1-14}
     \colorbox{CornflowerBlue}{\textbf{Ours}}  & 0.20 & 0.47 & 0.89 & 1.25 & 2.70 & 0.16 & 3.53 & 4.61 & 2.49 & 5.20 & 0.18 & 1.34 & 0.23 \\
    \bottomrule
    \end{tabular}
}
\vspace{1em}
\caption{Ablation results on Verse-Bench.}
\label{tab:ablation}
\vspace{-1em}
\end{table}